\section{Self-Modifying Software: de la ingeniería de prompts a la ingeniería de sistemas}

La característica más reconocible de OpenClaw es que «puede automodificarse». Esto no significa que el modelo haya adquirido conciencia de pronto, sino que la capa de sistema le permite al agente reescribir sus propias reglas de comportamiento dentro de ciertos límites.

\subsection{Qué diferencia hay entre «decirle qué no le satisface» y «que él modifique el sistema»}

En la IA conversacional tradicional, la retroalimentación del usuario normalmente solo afecta el siguiente turno de texto. La ruta de OpenClaw es distinta: la retroalimentación puede desencadenar cambios en la configuración, la documentación, los scripts y las cadenas de herramientas.

\begin{importantbox}{Expresión de ingeniería de la ruta de automodificación}
La entrevista puede abstraerse en un mecanismo de cuatro pasos:
\begin{enumerate}
    \item la retroalimentación del usuario se convierte en una intención ejecutable;
    \item el agente localiza el código o la configuración correspondiente;
    \item mediante la cadena de herramientas se envía y se valida el cambio;
    \item el resultado del cambio se escribe de vuelta en la memoria y en los archivos de reglas.
\end{enumerate}
Solo cuando estos cuatro pasos se hacen bien, el sistema pasa de «saber conversar» a «saber evolucionar».
\end{importantbox}

\subsection{Por qué este mecanismo es a la vez poderoso y peligroso}

\begin{knowledgebox}{Origen de su poder}
\begin{itemize}
    \item el ciclo de retroalimentación se acorta: el usuario no necesita esperar una iteración manual del producto;
    \item las optimizaciones locales se acumulan: el sistema se ajusta gradualmente al flujo de trabajo de cada persona;
    \item aumenta la reutilización: un proceso ya corregido puede trasladarse a tareas similares.
\end{itemize}
\end{knowledgebox}

El peligro es igual de evidente: una retroalimentación equivocada puede quedar «institucionalizada» dentro del sistema, y un atacante también puede aprovechar el prompt injection para inducir la escritura de reglas de largo plazo.

\subsection{Resumen de la sección}
La esencia del self-modifying software es trasladar el «aprendizaje» de la capa de pesos a la capa de sistema. Esto aumenta la velocidad de evolución, pero también adelanta la presión de seguridad y gobernanza a la etapa de diseño del producto.

